\chapter{Capa de plataforma}
\label{sec:capa-platform}

\filepath{platform/} agrupa nueve crates Rust y un paquete
TypeScript. Tienen en común ser \emph{superficies versionables}: lo
que publican es contrato público visto por consumidores internos o
externos.

\section{\filepath{platform/api} — la fachada única}

\code{GraphHunterApi} es el objeto central. Todos los consumidores
(Tauri, CLI, MCP, cualquier proceso embedder) obtienen una
instancia de esta struct y llaman métodos organizados por dominio:

\begin{itemize}
    \item \code{hunt::run_hunt}, \code{hunt::expand_node}
    \item \code{ingest::ingest_events}, \code{ingest::stream}
    \item \code{export::to_ocsf}, \code{export::to_iocs}
    \item \code{session::create}, \code{session::load},
    \code{session::save}
    \item \code{sentinel::connect}, \code{sentinel::stream}
    \item \code{agentic::run}, \code{agentic::review},
    \code{agentic::drift}
    \item \code{mapping_library::save},
    \code{mapping_library::apply}
    \item \code{dlq::list}, \code{dlq::retry}
    \item \code{invariants::check}
\end{itemize}

Son 55 endpoints HTTP cuando el Tauri los expone, y 54 tools MCP
cuando se exponen vía ese transporte. Pero la \emph{lógica} vive
una sola vez, en \filepath{platform/api}. ADR-001 formaliza esta
regla: la API es canónica; ninguna otra capa reimplementa parcial
ni ``mejora'' por dentro.

Esta decisión rompió el ciclo lógico que existía antes del
rediseño (\filepath{platform/api} $\to$ \filepath{apps/cli} $\to$
\filepath{core}). Hoy \filepath{apps/cli} es un cliente más de la
API, igual que Tauri.

\section{\filepath{platform/dsl}}

El DSL de hipótesis. La gramática actual acepta cadenas como:

\begin{lstlisting}[basicstyle=\ttfamily\small]
IP -[Connect]-> Host -[Auth]-> User {k=3}
\end{lstlisting}

La gramática vive en \filepath{grammar/v1.rs}. Un trait
\code{DslExtension} permite registrar nuevos tipos de entidad,
predicados y operadores sin modificar el parser base. ADR-005
detalla el esquema de versionado.

\section{\filepath{platform/parsers}}

Hogar del trait \code{LogParser} con siete implementaciones:
sysmon, sentinel, generic CSV, FortiAnalyzer, Cognito, IIS y un
parser fallback. La decisión de qué parser usar se toma en
\filepath{platform/api::format_registry}.

\section{\filepath{platform/sources}}

Trait \code{LogSource} para fuentes en vivo: Sentinel, Defender
XDR, Data Lake, y archivos locales. Es la contraparte ``pull'' de
\filepath{platform/parsers} (que trabaja sobre bytes) y vive
separada para poder reemplazarla sin tocar el parseo.

\section{\filepath{platform/canonical}}

Proyección OCSF v1.4 del modelo interno. Es la forma canónica que
se exporta hacia SIEMs externos o hacia el IOC exporter. La
dualidad ``modelo interno rico / modelo externo canónico''
permite evolucionar el interno sin cambiar el contrato público.

\section{\filepath{platform/vrl}}

Compilador de \code{FieldConfig} a VRL (Vector Remap Language). Se
usa en el camino de ingesta cuando la fuente viene con un mapeo
declarativo de campos. Contrato estable hacia clientes que comparten
configuraciones (ver \filepath{platform/api::mapping_library}).

\section{\filepath{platform/constrained-decode}}

Guardrail para salidas de LLM. Impone una gramática a la generación
(JSON strict, DSL de hipótesis, canonical query) para evitar que el
modelo produzca texto inválido. Vive aquí porque es una superficie
pública vista por el sistema agentic.

\section{\filepath{platform/siem}}

Streaming desde Sentinel y Elastic. Se extrajo de \filepath{apps/cli}
para romper el ciclo que describe ADR-001. Hoy lo consumen tanto
\filepath{platform/api} (para hacer ingest continuo desde el Tauri)
como \filepath{apps/cli} (para ingesta batch desde terminal).

\section{\filepath{platform/local-llm}}

Trait \code{LocalLlm} con backends \code{mock} y \code{candle}.
Permite al sistema agentic correr sin dependencias externas. Es
superficie extensible: agregar un backend nuevo es implementar el
trait, no tocar a los consumidores.

\section{\filepath{platform/mcp} (TypeScript)}

Servidor MCP con 54 tools organizadas en 10 dominios. No compila
junto a los crates Rust: tiene su propio \filepath{package.json} y
build (\code{npm run build}). Cada tool es un cliente HTTP liviano
contra \filepath{platform/api}. ADR-004 formaliza el contrato
\code{Tool<I, O>} y el versionado por tool.

\section{Qué tienen en común}

Los nueve crates Rust de \filepath{platform/} comparten tres
rasgos:

\begin{itemize}
    \item \textbf{Traits como contrato}: cada crate publica al menos
    un trait que define su extensibilidad (\code{LogParser},
    \code{LogSource}, \code{DslExtension}, \code{LocalLlm},
    \code{Tool} en el caso MCP).
    \item \textbf{DTOs versionados}: los tipos que cruzan fronteras
    (entre procesos o entre la API y el exterior) viven en
    \filepath{platform/api/src/dto/v1/} con compatibilidad semver.
    \item \textbf{Cero dependencia sobre aplicaciones}: un crate de
    plataforma nunca pide algo de \filepath{apps/}.
\end{itemize}

Esta uniformidad es lo que hace que el sistema se extienda sin
regresiones: un nuevo parser se añade en \filepath{platform/parsers}
sin que eso fuerce tocar ni el núcleo ni las aplicaciones.
